\documentclass[../../../include/open-logic-section]{subfiles}

\begin{document}
\olfileid{sth}{ord-arithmetic}{intro}

\olsection{పరిచయం}

\olref[ordinals][]{chap}లో క్రమసంఖ్యల సిద్ధాంతాన్ని అభివృద్ధి చేశాం. \olref[spine][]{chap}లో మిగిలిన స్థాయిక్రమం నిర్మించబడే వెన్నెముకగా క్రమసంఖ్యలను భావించవచ్చని చూశాం. అయితే క్రమసంఖ్యలకు ఆ పాత్ర ఒక్కటే కాదు. క్రమసంఖ్య అంకగణితం చేయడమూ మరో పని.

\olref[ordinals][idea]{sec}లో $\omega$, $\omega+1$, $\omega+\omega$ గురించి మాట్లాడినప్పుడు దీనిని ఇప్పటికే సూచించాం. అప్పుడు అనౌపచారికంగా మాట్లాడాం; ఇప్పుడు దాన్ని సరిగ్గా వివరించే సమయం వచ్చింది. అయితే ఈ అధ్యాయంలో తత్త్వశాస్త్రం ఎక్కువగా లేదని చెప్పాలి; సాంకేతిక అభివృద్ధులు మాత్రమే, అలాగే ఒకే పనిని రెండు భిన్నమైన మార్గాల్లో చేయవచ్చనే (కొద్దిగా) ఆసక్తికరమైన పరిశీలన ఉన్నాయి.

\end{document}
